\section{From task coordinates to virtual voltage inputs}
\label{sec:task-control}

Consider the electrical current dynamics already used in Part~III,
\[
 \dot{\vect i}=\vect f(\vect i,\omega)+\matr B(\vect i)\vect u,
 \qquad \vect y=\vect h(\vect i).
\]
The chain rule gives the task dynamics
\begin{equation}
 \dot{\vect y}=
 \underbrace{\gls{sym:jtask}\vect f}_{\text{task drift}}+
 \underbrace{\gls{sym:jtask}\matr B}_{\gls{sym:decoupling}}\vect u.
 \label{eq:task-output-dynamics}
\end{equation}
For outputs that depend directly on current, every regular channel has relative
degree one. If the decoupling matrix \(\matr A_h=\matr J_h\matr B\) is square
and nonsingular, the virtual task-rate input
\(\gls{sym:vtask}=\dot{\vect y}\) yields
\begin{equation}
 \boxed{\vect u=\matr A_h^{-1}
  (\vect v_{\mathrm{task}}-\matr J_h\vect f).}
 \label{eq:feedback-linearizing-input}
\end{equation}
This is input--output feedback linearization after, not before, its physical
meaning is clear: cancel the predicted task drift and choose converter
voltages that generate the desired task velocity. PMSM torque/flux control
has used precisely this idea \cite{lascu2012feedback,choi2016feedback}.

Invertibility is only the first test. The commanded \(\vect u\) must remain in
the stator inverter hexagon and field-converter interval. Close to a singular
\(\matr A_h\), small task commands demand large voltage and amplify model
error. Saturation breaks exact cancellation, so anti-windup or a constrained
least-squares allocation
\[
 \min_{\vect u\in\mathcal U}
 \|\matr A_h\vect u-(\vect v_{\mathrm{task}}-\matr J_h\vect f)\|_W^2
\]
is more honest than an unconstrained inverse.

\subsection{Three EESM states are a resource, not three equal obligations}

With independent stator and field converters, a linear EESM has three
electrical states and three input voltages. Candidate output triples include
\((M,\psi_s,\sigma)\), \((M,\psi_a,\alpha)\), and
\((M,\text{loss allocation},\text{voltage allocation})\). Each must pass the
rank test on \(\matr J_h\matr B\); a meaningful label is not enough.

A useful hierarchy is:
\begin{align*}
 \text{primary:}&\quad M\to M^\star,\\
 \text{secondary:}&\quad \text{minimize loss or regulate flux},\\
 \text{tertiary:}&\quad \text{voltage margin or thermal sharing}.
\end{align*}
Locally, \(\nabla M^{\mathsf T}\dot{\vect i}=0\) defines a two-dimensional
EESM torque-null tangent plane. One tangent direction can redistribute
stator/field excitation; the other can manage voltage. A two-current PSM has
only one such tangent direction. The existing \(z_0\) is one exact null
direction of the linear torque form, not the entire tangent plane at a general
point.

\begin{figure}[tb]
\centering
\resizebox{\linewidth}{!}{%
\begin{tikzpicture}[>=Latex,font=\scriptsize,node distance=5mm and 5mm]
\node[draw,rounded corners,fill=MidnightBlue!7] (t1) {\(M^\star,\omega\)};
\node[draw,rounded corners,right=of t1] (map) {MTPA/FW maps};
\node[draw,rounded corners,right=of map] (pi) {\(2\!+\!1\) current PIs};
\node[draw,rounded corners,right=of pi] (dec) {\(dq\) decoupling};
\node[draw,rounded corners,right=of dec] (pwm) {PWM/converters};
\draw[->] (t1)--(map);\draw[->] (map)--(pi);\draw[->] (pi)--(dec);\draw[->] (dec)--(pwm);
\node[left=3mm of t1,font=\bfseries] {CVC};

\node[draw,rounded corners,fill=ForestGreen!9,below=12mm of t1] (t2)
 {\(M^\star,\psi^\star,\) priorities};
\node[draw,rounded corners,right=of t2] (model) {task model/observer};
\node[draw,rounded corners,right=of model] (tc) {task-rate controller};
\node[draw,rounded corners,right=of tc] (alloc) {constrained input allocation};
\node[draw,rounded corners,right=of alloc] (pwm2) {PWM/converters};
\draw[->] (t2)--(model);\draw[->] (model)--(tc);\draw[->] (tc)--(alloc);\draw[->] (alloc)--(pwm2);
\node[left=3mm of t2,font=\bfseries] {task};

\draw[dashed,Orange,->] (pwm2.south) to[bend right=15] node[below]{measured currents, estimated flux} (model.south);
\end{tikzpicture}%
}
\caption{Conventional current-vector and task-space control stacks. Fewer
reference maps can be exchanged for observers, Jacobians, and constrained
input allocation.}
\label{fig:control-stacks}
\end{figure}


\subsection{Current references are optional, physical currents are not}

Conventional \gls{cvc} computes
\(M^\star\to\vect i^\star\to\vect u^\star\). DFVC, direct torque control,
deadbeat flux/torque control, and feedback linearization may instead compute
\((M^\star,\psi^\star)\to\vect u^\star\) directly. The recent EESM deadbeat and
unified torque-controller studies show how torque modules can be composed with
direct-flux dynamics \cite{rubino2026deadbeat,ionta2026unified}. Skipping an
explicit current reference can remove one map and one set of loops; it cannot
skip current and voltage constraints, state estimation, or the machine model
needed to predict the task response.

For the dynamic A-to-B problem, \((i_T,\alpha,z_0)\) suggests a transparent
schedule: change \(i_T\) rapidly with fast stator authority, move \(\alpha\)
more slowly to rebalance torque factors, and relax \(z_0\) along the remaining
null direction. This is a useful warm-start structure for the constrained
trajectory optimizer of Part~III, not proof that independently closed loops
produce the globally optimal path.
